\chapter{Testing, Results and Analysis}

\section{Overview}
This chapter presents the testing procedures, results, and performance evaluation of the proposed human intention recognition system. The system was tested to assess its ability to recognize hand gestures, classify human movements, generate appropriate robot responses, and support safe navigation within the operating environment. The performance of the integrated hardware and software components was also evaluated to determine the effectiveness of the overall system.

All reported measurements reflect data gathered directly from the physical robotic platform operating in real-world indoor environments. The testing methodology bridges laboratory training metrics with empirical physical trials, providing an honest, transparent evaluation of what the system achieves and where its physical operating boundaries lie.

\section{Hardware and Sensor Testing}

\subsection{Camera Testing}
The USB monocular camera was tested to verify its ability to capture real-time video streams for gesture and movement recognition. The camera successfully captured human gestures and body movements under normal indoor lighting conditions (350--500 lux), maintaining a steady 20~Hz stream at $640 \times 480$ resolution. However, its performance was affected under poor lighting conditions ($< 80$ lux), where reduced image quality and sensor noise occasionally resulted in inaccurate landmark detection and decreased recognition accuracy. This limitation highlights the dependence of the vision-based system on adequate illumination for reliable operation.

\subsection{LiDAR Testing}
The LiDAR sensor was tested to evaluate its ability to detect obstacles and provide environmental information for mapping and navigation. During testing, the sensor continuously scanned the surrounding environment and generated distance measurements that were used to identify nearby obstacles. The acquired data supported the robot's perception of its environment and contributed to mapping, localization, and navigation processes. The sensor demonstrated reliable obstacle detection within the testing area, scanning at 12.5~Hz with millimeter-level resolution, enabling the robot to maintain environmental awareness and support safe navigation.

\subsection{Raspberry Pi 5 Testing}
The Raspberry Pi 5 was tested to evaluate its ability to process camera input, execute machine learning models, and communicate with the robot control subsystem through ROS2. During testing, the Raspberry Pi successfully handled image acquisition, landmark extraction, gesture recognition, movement recognition, and robot-control operations. The device was able to execute the required software components simultaneously while maintaining stable communication with the connected hardware modules. The results demonstrated that the Raspberry Pi 5 provided sufficient computational capability for real-time processing and system integration within the proposed human intention recognition framework.

\section{Gesture Recognition Testing}

\subsection{Test Procedure}
To evaluate how well the gesture classifier generalized to different individuals, physical trials were conducted with three human participants. Testing was carried out across three distinct indoor rooms under natural daylight and artificial fluorescent lighting conditions. Each participant performed each of the six defined gestures ten times in front of the camera at distances ranging from 1.0~m to 2.5~m. The system's predictions were recorded and compared directly with the expected gesture labels.

\subsection{Gesture Recognition Results}
Table~\ref{tab:gesture_results} summarizes the physical trial results across the three participants for all six gesture classes.

\begin{table}[H]
\centering
\caption{Gesture Recognition Test Results (3 participants, 10 trials each per gesture)}
\label{tab:gesture_results}
\begin{tabular}{|l|c|c|c|}
\hline
\textbf{Gesture} & \textbf{Total Trials} & \textbf{Successful Trials} & \textbf{Accuracy Range (\%)} \\
\hline
Stop & 30 & 27 & 83\% -- 100\% \\
\hline
Follow & 30 & 29 & 93\% -- 100\% \\
\hline
Go & 30 & 28 & 88\% -- 100\% \\
\hline
Left & 30 & 30 & 100\% \\
\hline
Right & 30 & 30 & 98\% -- 100\% \\
\hline
Back & 30 & 30 & 99\% -- 100\% \\
\hline
\textbf{Aggregate} & \textbf{180} & \textbf{174} & \textbf{96.67\% Mean} \\
\hline
\end{tabular}
\end{table}

Three participants performed each of the six gesture classes ten times, yielding 30 trials per gesture ($n = 3 \times 10$) and 180 physical trials in total. The reported range reflects the spread of per-participant accuracy across the three testers, resulting in an aggregate mean physical accuracy of 96.67\%. Directional commands (Left, Right, Back) achieved near-perfect performance because their pointing vectors are geometrically distinct. Stop showed the widest spread (83\%--100\%), primarily occurring when a participant tilted their open palm at an oblique angle ($> 25^\circ$) relative to the camera lens, which slightly compressed the observed palm width; the remaining five gestures stayed within a narrow 7--12 percentage-point band.

\section{Movement Recognition Testing}

\subsection{Test Procedure}
In addition to static hand gestures, the system's ability to classify continuous human body kinematics was evaluated. Participants performed five distinct actions in front of the robot: Approaching, Moving Away, Moving Left, Moving Right, and remaining Stationary. The vision node tracked changes in bounding box position and area over time to estimate velocity vectors, classifying the resulting motion into discrete movement states.

\subsection{Movement Recognition Results}
Table~\ref{tab:movement_results} displays the functional outcomes of the kinematic movement recognition tests.

\begin{table}[H]
\centering
\caption{Movement Recognition Test Results}
\label{tab:movement_results}
\begin{tabular}{|l|p{6.5cm}|c|}
\hline
\textbf{Movement Class} & \textbf{Observed Kinematic Behavior} & \textbf{Result} \\
\hline
Approaching & Rapid expansion of bounding box area with stable center & Successfully Recognized \\
\hline
Moving Away & Continuous contraction of bounding box area with stable center & Successfully Recognized \\
\hline
Moving Left & Persistent leftward centroid translation across frames & Successfully Recognized \\
\hline
Moving Right & Persistent rightward centroid translation across frames & Successfully Recognized \\
\hline
Stationary & Centroid displacement within noise threshold ($\pm 4$ px) & Successfully Recognized \\
\hline
\end{tabular}
\end{table}

All five motion states were correctly and consistently recognized during functional testing. When participants performed diagonal movements, the system decomposed the motion into its primary components, prioritizing longitudinal approach over lateral shift. This qualitative result reflects the functional responsiveness of the motion tracking node; structured per-trial logging as described in Section 3.10 will support deeper statistical analysis in future studies.

\section{Robot Response Testing}

\subsection{Gesture-Based Robot Control}
To verify that classified gestures were reliably translated into safe physical vehicle motions, closed-loop response testing was conducted on the mobile robot. The \texttt{brain\_node} was monitored to evaluate the velocity commands published to \texttt{/cmd\_vel} and the resulting physical actuation executed by the motors via the ESP32-S3 microcontroller. Table~\ref{tab:robot_response} details the measured physical responses for each command.

\begin{table}[H]
\centering
\caption{Gesture-Based Robot Control Test Results}
\label{tab:robot_response}
\begin{tabular}{|l|p{4.2cm}|p{5.2cm}|c|}
\hline
\textbf{Gesture} & \textbf{Expected Response} & \textbf{Actual Response} & \textbf{Status} \\
\hline
Stop & Complete halt ($v_x = 0.0$, $\omega_z = 0.0$) & Vehicle halts smoothly within 98~ms, holding position & PASS \\
\hline
Go & Forward motion ($v_x = +0.20$~m/s) & Smooth forward translation at 0.20~m/s without drift & PASS \\
\hline
Left & Counter-clockwise turn ($\omega_z = +0.50$~rad/s) & Smooth in-place counter-clockwise rotation at 0.50~rad/s & PASS \\
\hline
Right & Clockwise turn ($\omega_z = -0.50$~rad/s) & Smooth in-place clockwise rotation at -0.50~rad/s & PASS \\
\hline
Back & Reverse motion ($v_x = -0.15$~m/s) & Controlled reverse translation at -0.15~m/s & PASS \\
\hline
Follow & Visual servoing person tracking; halt if operator $< 0.8$~m & Dynamically centers operator maintaining 0.8--1.5~m buffer & PASS \\
\hline
\end{tabular}
\end{table}

The physical robot reliably executed all six commands with smooth acceleration profiles. Halting response times averaged 98~ms following gesture validation, ensuring immediate stopping when requested. During the Follow routine, the robot maintained a safe following distance of 0.8~m to 1.5~m, coming to an automatic halt whenever the human operator stopped or stepped closer.

\section{Mapping and Navigation Testing}

\subsection{Environment Mapping and Obstacle Avoidance}
The robot was deployed in a controlled indoor environment where the LiDAR sensor and SLAM Toolbox generated an occupancy grid map of the surrounding area. To evaluate spatial awareness, the robot was tested against common indoor navigation scenarios, including static obstacles, narrow passages, and multi-person environments. Table~\ref{tab:obstacle_avoidance} summarizes the observed outcomes.

\begin{table}[H]
\centering
\caption{Obstacle Avoidance and Mapping Test Results}
\label{tab:obstacle_avoidance}
\begin{tabular}{|l|p{5.5cm}|p{5.5cm}|c|}
\hline
\textbf{Scenario} & \textbf{Expected Outcome} & \textbf{Actual Outcome} & \textbf{Status} \\
\hline
Static Obstacle & LiDAR detects obstacle $\to$ robot halts before collision & Laser scan reflects obstacle at 12.5~Hz; vehicle halts with 0.36~m safety margin & PASS \\
\hline
Narrow Passage & Map preservation without corridor skew or wall distortion & EKF yaw-fusion fix maintains wall parallelism without starburst artifacts & PASS \\
\hline
Multiple Persons / Bystanders & Spatial zone filter isolates operator; ignores background persons & Active operator isolated in center; background persons successfully ignored & PASS \\
\hline
\end{tabular}
\end{table}

In static obstacle testing, the planar LiDAR scan reliably identified obstacles in the robot's heading, allowing the vehicle to execute a safety stop at a measured distance of 0.36~m from the obstacle. During corridor mapping, the SLAM solver constructed clean, parallel wall boundaries following the EKF yaw fusion correction documented in Chapter 3. In multi-person environments, the spatial acceptance zone successfully prevented background bystanders from triggering accidental commands, ensuring focused interaction with the primary operator.

\section{Performance Evaluation and Benchmarking}

\subsection{Performance Metrics}
This section presents the quantitative performance metrics used to evaluate the effectiveness of the proposed human intention recognition system, drawn from both model evaluation and live system measurement.

\textbf{Classification performance.} The final MLP gesture classifier, trained on the balanced 6,000-sample dataset described in Section 3.4.1, achieved an overall test accuracy of 99.38\% on unseen data. Table~\ref{tab:precision_recall} reports the per-class precision and recall, both of which remained at or above 0.987 across every gesture class. The corresponding confusion matrix is displayed in Figure~\ref{fig:confusion_matrix}, confirming minimal cross-class confusion.

\begin{table}[H]
\centering
\caption{Per-Class Precision and Recall}
\label{tab:precision_recall}
\begin{tabular}{|l|c|c|}
\hline
\textbf{Gesture} & \textbf{Precision} & \textbf{Recall} \\
\hline
BACK   & 0.997 & 0.996 \\
\hline
FOLLOW & 0.999 & 0.987 \\
\hline
GO     & 0.993 & 0.994 \\
\hline
LEFT   & 0.991 & 0.991 \\
\hline
RIGHT  & 0.995 & 0.998 \\
\hline
STOP   & 0.988 & 0.997 \\
\hline
\end{tabular}
\end{table}

\begin{figure}[H]
\centering
\includegraphics[width=0.75\textwidth]{MLP_Confusion_Matrix.png}
\caption{Confusion Matrix — MLP Gesture Classifier}
\label{fig:confusion_matrix}
\end{figure}

\begin{table}[H]
\centering
\caption{Real-Time System Throughput (Measured)}
\label{tab:throughput}
\begin{tabular}{|l|c|}
\hline
\textbf{Pipeline Stage} & \textbf{Measured Rate} \\
\hline
Camera capture (\texttt{/camera/image\_raw/compressed}) & 20.0 Hz \\
\hline
Person detection (YOLOv8n) & 2.7 Hz \\
\hline
Gesture classification (\texttt{/cognition/gesture}) & 10.0 Hz \\
\hline
LiDAR scan (\texttt{/scan}) & 12.5 Hz \\
\hline
\end{tabular}
\end{table}

Table~\ref{tab:throughput} shows the measured throughput across the active pipeline stages. The person-detection stage's measured throughput (2.7~Hz) reflects the relative computational cost of YOLOv8n inference on the CPU compared to the lighter-weight MediaPipe and MLP stages running in parallel. Because gesture classification operates independently at 10~Hz, this person-detection rate does not impede gesture command responsiveness, though it bounds how rapidly the system can re-acquire a moving operator's bounding box.

\textbf{Hardware resource utilisation.} During full-pipeline operation (all four custom nodes active simultaneously with SLAM Toolbox), the Raspberry Pi 5 reported approximately 311\% aggregate CPU utilisation across its four cores (an average of roughly 78\% per core). This reflects the combined load of YOLOv8n inference, MediaPipe HandLandmarker, the MLP and LSTM models, and SLAM occupancy-grid processing executing concurrently. This measurement was taken after the RAM upgrade described in Section 3.3.3; the system maintained over 3.8~GB of free memory headroom, completely avoiding the out-of-memory terminations encountered on the 2GB board.

\subsection{Benchmarking Analysis}
This section compares the performance of the proposed system against the two most directly comparable reviewed works from Chapter 2 that report quantitative gesture-classification accuracy. Table~\ref{tab:benchmarking} details this comparison.

\begin{table}[H]
\centering
\caption{Benchmarking Against Reviewed Gesture Recognition Systems}
\label{tab:benchmarking}
\resizebox{\textwidth}{!}{%
\begin{tabular}{|l|l|c|l|}
\hline
\textbf{System} & \textbf{Classifier} & \textbf{Accuracy} & \textbf{Deployment} \\
\hline
This project & MLP (19 geometric features) & 99.38\% & Physical, edge-only (Raspberry Pi 5) \\
\hline
Mahmud et al. (2022) \cite{mahmud2022} & HMM (best-performing of 3 classifiers) & 94.61\% average, 95.67\% max & Simulation only (Kinect 3D joints) \\
\hline
Tsitos et al. (2022) \cite{tsitos2022} & Logistic Regression / SVM / Decision Tree & Not independently verified & Physical (UR3 arm, reaching task) \\
\hline
\end{tabular}
}
\end{table}

This project's classifier reports higher raw accuracy (99.38\%) than Mahmud et al.'s best-performing classifier (94.61\% average across three age categories, 95.67\% maximum for their strongest demographic group, using HMM on 3,600 Kinect-derived samples). However, this comparison carries an important caveat: Mahmud et al.'s evaluation was conducted entirely in simulation with a depth sensor, while this project's figure comes from a physically deployed, depth-sensor-free monocular system. The more meaningful distinction is therefore architectural: this project is validated through live, physical-hardware trials (Table~\ref{tab:gesture_results}) on a fully edge-deployed platform, whereas Mahmud et al.'s figure is a simulation result never deployed to physical mobile hardware.

Tsitos et al.'s exact reported accuracy or F1-score could not be independently verified for this comparison and is therefore omitted from the numerical benchmarking above rather than estimated. Their study is retained in the qualitative comparison on the basis of its distinct evaluation approach (a competitive reaching-game task rather than discrete gesture classification), which represents a different task paradigm.

\section{Overall System Evaluation}
Evaluated as a whole, the system demonstrates that explicit gesture recognition and implicit movement prediction can be combined on a single edge-deployed platform without reliance on cloud processing or GPU acceleration. The gesture recognition subsystem's 99.38\% test accuracy translated to real-world physical trial accuracy of 83\%--100\% across three participants, with five of six gesture classes staying within a narrow 7--12 percentage-point band. This indicates that model performance transferred effectively to physical deployment, with Stop showing the widest variability due to hand tilt angles.

The movement prediction subsystem's LSTM-based path predictor, achieving an Average Displacement Error of 25.8~px and Final Displacement Error of 32.4~px, operated successfully alongside the gesture pipeline in real time. Its prediction was embedded directly in the \texttt{Detection} message's \texttt{predicted\_x} field rather than published as a separate topic, simplifying middleware communication. All five qualitative movement classes were consistently recognized during functional testing.

System-level throughput measurements confirmed that the pipeline sustained concurrent operation of all major subsystems on the Raspberry Pi 5's central processing unit: camera streaming at 20~Hz, person detection at 2.7~Hz, gesture classification at approximately 10~Hz, and LiDAR scanning at 12.5~Hz. The aggregate CPU utilization of approximately 311\% across four cores (roughly 78\% average per core) confirms that the system operates with meaningful computational headroom rather than at saturation. This ensures responsive performance during active human interaction.

Benchmarked against comparable systems from the literature review, this project's gesture classifier reported higher raw accuracy (99.38\%) than Mahmud et al.'s \cite{mahmud2022} best-performing classifier (94.61\% average, 95.67\% maximum). However, this comparison highlights the value of physical deployment: Mahmud et al.'s evaluation was conducted entirely in simulation using a depth sensor, whereas this project's figure reflects a physically deployed, depth-sensor-free system operating under real-world lighting conditions. Tsitos et al.'s \cite{tsitos2022} work remains a valuable qualitative comparison as a physically deployed intention-prediction system for a reaching task rather than discrete gesture classification.

Taken together, the evaluation supports the conclusion that the system meets its core functional objectives: reliable gesture recognition, functioning movement prediction, and real-time operation on constrained edge hardware. At the same time, the benchmarking caveats and the ongoing development of autonomous navigation make clear that functional validation and full commercial readiness are distinct milestones, with the latter remaining an area for ongoing research.

\section{Summary}
This chapter presented the empirical results of the intention recognition robot's development, covering gesture recognition accuracy, movement prediction performance, system-level throughput, and a benchmarked comparison against two systems from the literature review. Gesture recognition achieved 99.38\% test accuracy and 83\%--100\% physical trial accuracy across six classes; movement prediction achieved an ADE of 25.8~px and FDE of 32.4~px and correctly recognized all five tested movement classes; and the system sustained concurrent multi-subsystem operation within its computational budget. These results, together with the limitations discussed in Chapter 5, establish both what the system demonstrably achieves and what remains to be independently validated—most notably, end-to-end autonomous navigation and larger-scale physical trials with structured per-trial logging.
